% !TeX root = ../main.tex

\xchapter{结论与展望}{Conclusion and Future Work}

\xsection{工作总结}{Work Summary}

本论文工作依托于教育部基础学科和交叉学科突破计划(JYB2025XDXM116)及国家自然科学基金(62477036)。针对平面几何长链条推理中面临的长步骤评测基准缺失、多模态感知误差累积以及逻辑发散等关键挑战，系统地展开了数据构建、算法设计与系统实践。具体研究内容包括以下三个方面：

第一，针对长链条推理评测基准与形式化基础缺失问题，构建了一个面向长步骤推理的多模态平面几何基准 GeoLStep。该工作从真实教育场景出发，收集并标注了高难度的几何综合题，设计了将自然语言与几何配图通向严格符号逻辑的形式化转换链路，并引入了答案准确率与过程逻辑正确率的双重评估机制。该基准不仅填补了复杂逻辑评测的空白，也为后续过程级强化与自动化验证提供了坚实的数据底座。

第二，针对跨模态长链条推理中的逻辑失稳与视觉干扰问题，提出了一种融合神经感知与弱符号规则反馈的两阶段增强方法。第一阶段在纯文本环境下，通过组相对策略优化强化模型的基础演绎范式；随后，第二阶段在多模态环境下引入弱符号验证器，对自回归生成的中间步骤进行实时合法性检验与冗余判定。该方法将步骤级反馈转化为策略优化的局部奖励信号，有效规范了大模型的生成轨迹，显著抑制了长链推演后期的误差发散现象。

第三，为了验证理论方法在真实教育场景下的工程可行性与应用落地，设计并实现了一个交互式的平面几何问题求解原型系统。该系统集成了多模态数据输入、多异构模型并发调度与对比，以及推理步骤可视化展示等功能，支持多模型切换对比。通过前后端分离架构与异步并发调度机制，满足了自动解题在多并发场景下的扩展需求。

本研究以“神经感知与规则约束协同”的策略为核心，提出了一套面向长链条几何推理的评测与求解方法，同时通过系统平台验证了方法的适用性与可操作性。为复杂多模态逻辑推理任务提供了新的技术路径与系统解决方案，推动了智能解题技术在教育数字化中的深度融合。

\xsection{工作展望}{Future Work}

尽管本研究在长链条平面几何推理的基准构建与强化求解方面取得了一定的研究成果，并成功构建了具备多模型对比与可视化能力的系统平台，验证了提出方法的有效性，但仍存在一些局限性，未来仍有多个方向值得进一步探索与深入优化。主要展望如下：

第一，增强对隐性空间拓扑关系的联合感知能力。本研究虽然通过形式化转换提取了关键几何要素，但在实际求解中，模型对“点在线上的相对位置”等隐性空间拓扑特征仍缺乏足够的约束。未来的研究可尝试引入图神经网络深入挖掘几何图元间的连续空间分布特征，实现显式拓扑先验与隐式神经表征的深度融合，从而提升模型在复杂几何图形下的视觉感知鲁棒性与空间推理精度。

第二，优化弱符号验证器的动态反馈机制。本研究提出的两阶段强化方法主要依赖基于预定义规则的弱符号验证器提供步骤级反馈。未来可探索更灵活、动态的中间状态评估模型或强弱结合的神经符号协同机制。例如，引入外部定理搜索树或自适应的价值网络，使多模态大语言模型在生成阶段获得更精细、前瞻性的逻辑制导，进一步提升超长推理链条的正确率与泛化性。

第三，探索更复杂的动态交互场景与持续学习机制。本研究中构建的求解系统具备了基础的多模态接入与并发评测能力，但仍以离线模型调度为主。在实际的智能辅导应用中，解题过程往往是动态交互的。未来可进一步研究支持人机协同解题的动态交互框架，允许系统根据用户的局部提示或纠错实时调整推演路径。同时，探索在线学习机制，使模型在接收新题型的同时不断自我迭代，以更好服务于个性化导学等真实教学场景。

综上所述，未来的研究将从多模态拓扑感知、神经符号动态耦合、持续交互学习等多个方向持续推进，进一步完善几何智能推理的理论体系与工程落地能力，推动其在智慧教育与复杂逻辑推理场景中的智能化升级。